% !TeX root = ../papers/nc-p3.tex
\section{楕円五次窓の二捩れへの限定}
\label{sec:elliptic-quintic-cubic-overlaps}
\begin{theorem}\label{thm:elliptic-quintic-torsion}
前節の仮定で、各添字に独立な線型基底変更と点曲線の同時再パラメータ化を許すと、
楕円五次窓の候補は $j_{\mathrm i},j_{-\mathrm i}$ に対する各一つだけである。
具体的には $a=(1,0), b=(-1,0)$ とし、隣接辺に
\[
 R(a),R(b),R(a),R(b),\ldots,\qquad
 R(v)=\ker\bigl(V\otimes V\to H^0(L\,T_v^*L),\quad f\otimes g\mapsto f\,T_v^*g\bigr)
\]
を置く。ここで $V$ はそれぞれの明示楕円五次系、$T_v$ は平行移動である。
逆順の二捩れ列は同じ添字保持同型類を与える。両符号が実現するなら二類は相異なる。
\end{theorem}
\begin{proof}[三次重なりの指標分解]
複素共役で片方の符号を扱えば足りる。以下 $\epsilon=\mathrm i$、$H=\ker(2-\mathrm i)$ とする。
隣接二辺は $R(t),R(t+d)$、$d\in E[10]$ で、各辺の差は $H$ を法として選べる。
実際、頂点の平行移動に $H$ の元を独立に加えても、埋込みの線型座標変更で吸収できる。
二辺が有限Frobenius窓に属するなら
\[
 (R(t)\otimes V)\cap(V\otimes R(t+d))
\]
は四次元である。その四つの $H$ 指標は全て異なる。実際、全四位置の関係を保つ $H$ の作用は
一次元の四次交叉直線も保つ。この直線を第四因子で縮約して、第四文字の双対（共通指標ひねり）と
三次交叉との同型を得る。$V$ 自身の四指標は前節により相異なる。

\smallskip
\noindent\textit{良い還元における全幾何点の検算。}
曲線と十捩れを標数181、$\mathrm i=19$ に還元する。十捩れは有限エタールで、
全100点がこの素体上有理である。$H$ を法として二階差 $d$ は20通り。
\path{computations/verify_elliptic_quintic_overlaps.py}が行う厳密検算の内容を示す。
指標を $1,2,3,4\pmod5$ とするよう線型化し、その基底での点座標を $z(P)$ とする。
$h$ ブロックの未知係数を三指標の和が $h$ の語に限る（12または13列）。
$q$ は相異なる三つの $H$ 剰余類の代表を動く。左側は最初の二位置を
$z(q),z(q+t)$ で評価し、右側は後ろ二位置を $z(q),z(q+t+d)$ で評価する。
各々残る四つの添字で係数を取り、計24行とする。
これらの行の共通核は当該指標の全三次交叉である。
実際、対の評価で得る曲線切断は、残る一添字を固定すれば一定指標を持つ。
次数10の束は巡回五重商で次数2の束に降下し、一指標の切断が三剰余類で消えれば恒等零である。

一般階数の証明には各ブロックで $t=(6,36)$ において全列階数を持つ一小行列を選ぶ。
この最大次小行列は $t$ の関数であり、各行の極はその行の変数点が $O$ となる一点で高々5次である。
有限な極の $x$ 座標 $x_j$ と使用行数 $m_j$ に対して
$\prod (x-x_j)^{5m_j}$ を掛け、行列式を $A(x)+y B(x)$ に直す。
無限遠の極上限は $10\sum m_j+5m_O\le130$。よってそのRiemann--Roch基底
$x^i,x^i y$ に対し、正則な有理点での全列階数の補間行列で厳密に復元できる。
各ブロックのノルム $N_h=A^2-(x^3-x)B^2$ を作る。清算に用いた極でも、
射影的な小行列が零ならば清算後の値は零である（各行の極上限で既に清算している）。
無限遠は別に直接調べる。

四指標が落ちる点では任意の相異なる $h,j$ に対して $N_hN_j=0$。
全十積の最大公約元を各 $d$ で因数分解すると、全ての因子は一次で、
全ての根の上で曲線の $y$ も素体にある。従ってここでの候補点の直接階数検査は
代数閉包上の全候補を覆う。残る点を $(d,t)$ の $H$ 剰余類で示すと
\[
 (0,0),\qquad (k,a),\qquad(k,b),\qquad k=(0,0).
\]
三点とも、指標0は単射で、他の四指標の核は各一次元である。

\smallskip
\noindent\textit{標数零への移行。}
素点181での離散付値環を必要に応じ有限拡大し、各捩れ点はそのエタール持上げを用いる。
対象の階数条件は $E$ 上の閉条件である。曲線・捩れ・線型系が代数的に定義されているため、
特定の捩れ点以外に標数零の解があるなら、代数閉体 $\overline{\mathbb Q}$ 上にもある。
固有性でその解を付値環上に延長できる。特殊化で上の三剰余類のいずれかに入るので、
各類において捩れ点と同じところへ特殊化する別解を排除すればよい。

上の三点で、評価用 $q$ を極を避けて取り直し、局所変数 $z$ で評価行列を二次まで展開する。
不変微分に双対な $\partial x=2y,\ \partial y=3x^2-1$ による
$1,\partial,\partial^2/2$ の係数は、標数181でもこの位数まで正則局所座標を与える。
単元で消去し、列余階数が1なら、残った一列の各係数の最小消滅次数が最大次小行列の最小次数になる。
五指標での結果は順に
\[
 (0,2,2,2,2),\qquad (0,1,1,1,1),\qquad (0,1,1,1,1).
\]
$H$ による平行移動に関しては埋込みの基底変更で同じ局所次数となる。

$(k,a),(k,b)$ では標数零でも三次交叉が四次元である。
実際 $u=\pm1$ に対する二捩れ移動の座標は
$x'=u+(3u^2-1)/(x-u),\ y'=-(3u^2-1)y/(x-u)^2$ である。
二系の16積を $(x-u)^3$ で清算し $y^2=x^3-x$ を使って核を取ると、二辺の挿入行列の階数は
両順とも $\mathbb Q(\mathrm i)$ 上で44（同コード末尾）となる。
従って四指標の核は捩れ点の持上げで実際に残る。上の単純零の小行列は、
捩れ切断からの局所変数 $z$ で、完備局所環において $z$ と単元の積である。

$(0,0)$ では二辺とも同じ $R(t)$ で、$t=0$ で交換子六本となる。
各非零指標の一次元核は任意の一次変形で少なくとも一次まで持ち上がる。
実際、変形関係を $ab-ba-zP_{ab}+O(z^2)$ と取り、
三文字の交代syzygyを左右からこれで持ち上げた際の一次の差は
$\sum_{\rm cyc}(P_{ab}c-cP_{ab})$ である。この差は可換三次商で零なので
零次関係の左・右挿入で消せる。指標射影で持上げを同じ指標に保てる。
従って当該最大次小行列は捩れ切断で $z^2$ により整除される。
ここでの関係束は乗法が全射で局所自由であり、この議論は付値環上でも行える
（または標数零での定数・一次係数の消滅と付値環の単射を使う）。
特殊ファイバーでちょうど二次の小行列を選べば、完備局所環で $z^2$ と単元の積である。
以上、どの場合も同じ剰余点への他の持上げを許さない。評価点と指標基底は滑らかな曲線、
エタール捩れおよび相異なる五乗根を用いて持ち上げればよい。

隣接差0の剰余類は前節により排除済みである。従って各位置の差は $a$ と $b$ を $H$ を法として
交互に繰り返す。各埋込みの座標変更で剰余類の代表を順に正規化でき、
全辺で表示の列を得る。$[\mathrm i]$ は $a,b$ を交換し、どちらの $j_\epsilon$ にも
$[1+\epsilon x^2:-x:\mathrm i y:-\mathrm i xy]$ と線型に作用するので逆順は同型である。
最後に、三割線入射のうち第一射影次数1の成分が決める単数 $\beta$ は
曲線同型で共役される。楕円曲線の準同型環は可換で、平行移動は準同型部分を変えないため、
$\mathrm i,-\mathrm i$ は区別される。
\end{proof}
